\chapter{Tekniska motåtgärder mot social engineering}
Kapitlet kommer att analysera tekniska motåtgärder som kan användas för att förebygga, identifiera och begränsa social engineering-attacker. 
Som stöd används de angreppsmetoder och fallstudier som tidigare presenterats där en analys av hur tekniska säkerhetsåtgärder kan förstärka motståndskraften mot 
social manipulation. Inledningsvis diskuteras metoder för identifiering och detektion av social engineering-relaterade attacker. Därefter behandlas artificiell 
intelligens som en potentiell försvarsmekanism, med fokus på dess effektivitet och begränsningar. Avslutningsvis diskuteras autentisering och åtkomstkontroll som 
tyngdpunkt i skyddet mot obehörig åtkomst och identitetsbedrägeri.

\section{Identifiering och detektion av social engineering-attacker}
Eftersom social engineering riktar sig mot mänskligt beslutsfattande kan identifiering av dessa attacker vara särskilt svår. Till skillnad från traditionella intrång, 
där skadlig kod eller nätverksintrång kan upptäckas genom teknisk övervakning, sker social manipulation ofta genom kommunikationskanaler som e-post, 
short message service (sms) eller sociala medier \cite{se_threat_defense_survey, algarni2019se_survey}. Av denna anledning fokuserar en stor del av detektionsarbetet 
på tekniska filtrerings- och analysmekanismer. I nätfiskedetektion används exempelvis e-postfiltrering samt analys av URL-strukturer, domänvalidering och 
maskininlärningsbaserade klassificeringsmodeller för att identifiera suspekta meddelanden innan de når offret \cite{basnet2023anti_phishing}. Systemen kan exempelvis 
analysera avsändarinformation, språkmönster och beteendeavvikelser för att markera potentiellt skadligt innehåll.

Utöver de ovannämnda detektionssystemen används ofta säkerhetsinformation och händelsehantering (SIEM-system) för att samla in och korrelera säkerhetsloggar från olika 
delar av en organisations infrastruktur. Genom kombinationen av SIEM och intrångsdetekteringssystem (IDS) kan misstänkt aktivitet identifieras i realtid. 
Enligt \cite{siem_comprehensive_study_2022} visar integrationen av SIEM och IDS med maskininlärningsmodeller en förbättrad förmåga att upptäcka ovanliga inloggningar 
eller typiska åtkomstförsök. Moderna organisationer har även börjat implementera Security Orchestration, Automation and Response (SOAR-system). Ett SOAR-system integrerar 
flera olika säkerhetsverktyg i en gemensam miljö som möjliggör automatiserad hantering av säkerhetsincidenter \cite{almohri2024soar}. Genom insamlade data från exempelvis 
SIEM-system, e-postfilter och hotinformationsdatabaser kan SOAR automatiskt initiera åtgärder såsom blockering av konton, isolering av enheter eller återkallande av 
skadliga länkar. Därmed möjliggör SOAR snabbare identifiering och respons vid social engineering-relaterade incidenter.

Trots alla dessa system kvarstår utmaningar. Eftersom social engineering bygger på psykologi snarare än teknik, kan välkonstruerade attacker fortfarande undvika 
traditionella säkerhetsfilter \cite{algarni2019se_survey}. Därför betonas i organisatoriska sammanhang vikten av att kombinera tekniska detektionssystem med 
organisatoriska rutiner och användarmedvetenhet för att uppnå optimalt skydd.

\section{Artificiell intelligens som försvarsmekanism}
Utvecklingen av AI har dragit med sig både för- och nackdelar inom cybersäkerhet. I dagens cyberhotlandskap har AI sett en ökad användning av hotaktörer för att 
exempelvis skapa mer övertygande nätfiskeattacker eller djupfejkmaterial, medan samma teknologi har blivit ett centralt verktyg för att motverka dessa social 
engineering-attacker \cite{zhang2024digital_deception}.

Orsaken till att AI har blivit ett viktigt verktyg för att motverka social engineering-attacker är dess förmåga att analysera stora kvantiteter av data och identifiera mönster 
som kan vara svåra att upptäcka manuellt. I nätfiskedetektion används exempelvis maskininlärningsmodeller för att enkelt markera e-postmeddelanden baserat på språk, 
metadata, URL-strukturer och avsändarbeteende \cite{basnet2023anti_phishing}. Genom att kontinuerligt träna modellerna med nya data kan försvaret mot social engineering 
upprätthållas i realtid. Enligt \cite{zhang2024digital_deception, basnet2023anti_phishing} kan AI även användas för att förstärka användarstöd och beslutsfattande. Genom specialiserade varningssystem 
och kontextbaserade säkerhetsindikatorer kan AI minska risken för impulsivt agerande vid en social engineering-attack. 

Som redan nämnt fungerar AI inte endast som en universallösning. Eftersom hotaktörer också använder AI för att skapa mer sofistikerade angrepp, uppstår det en form av kapprustning \cite{zhang2024digital_deception}. 
För att hållas i kapp måste därför försvarssystem kontinuerligt uppdateras och kombineras med organisatoriska säkerhetsrutiner \cite{algarni2024comprehensive_taxonomy}.



\subsection{Användningen av maskininlärning för detektion}
Som nämnts ovan utgör maskininlärning (ML) en central del av AI-baserade försvarssystem mot social engineering. Till skillnad från traditionella försvarssystem, som kräver fördefinierade kriterier, 
kan användningen av övervakad och oövervakad maskininlärning användas för att identifiera komplexa mönster även i stora datamängder \cite{basnet2023anti_phishing}. 
Användningen av metoderna möjliggör anpassning till kontinuerligt förändrande attacker.

I nätfiske-relaterad detektion används ofta övervakad inlärning, där modeller tränas på specifika dataset bestående av skadliga och ofarliga meddelanden. Genom att fokusera på exempelvis 
analys av språkliga strukturer, avsändardomäner, metadata och förekomsten av språkliga indikatorer på brådska, såsom tidsbegränsningar, krav på omedelbar åtgärd eller hot om negativa konsekvenser, 
kan modellen klassificera nya meddelanden med större sannolikhet \cite{basnet2023anti_phishing}. Metoden möjliggör konstant förbättring av detektionsprecisionen i takt med nya data. Utöver övervakad 
inlärning används också oövervakad inlärning för att identifiera avvikande beteendemönster. I dessa fall analyseras en användares aktivitetsmönster, som exempelvis inloggningstider och geografisk placering. 
Om användarens beteende avviker i stor grad från det normala kan systemet generera en varning om ett potentiellt äventyrande konto \cite{siem_comprehensive_study_2022}. Metoden är därför särskilt användbar i 
situationer där hotaktören redan har fått åtkomst till autentiseringsuppgifter genom lyckad social engineering.  

ML används även i kombination med SIEM- och IDS-system för att möjliggöra realtidsanalys av nätverkstrafik och autentiseringsloggar \cite{siem_comprehensive_study_2022}. Genom att använda flera datakällor 
kan systemet skapa en helhetsbild av det potentiella hotet. Integrationen möjliggör förbättrad identifiering av komplexa attacker som kombinerar tekniskt utnyttjande med social manipulation. Oavsett sina styrkor 
har ML också begränsningar. Modellerna är mycket beroende av kvaliteten på träningsdata och kan vara särskilt känsliga för så kallade fientliga attacker, där hotaktörer medvetet manipulerar inputdata för att undvika 
detektion \cite{algarni2024comprehensive_taxonomy}. ML bör därför förstås som en del av en helhetsstrategi istället för en självständig lösning.

\subsection{Begränsningar och risker med AI}
Även om AI erbjuder avancerade möjligheter för detektion och respons med hjälp av ML, är teknologin ingen universallösning mot social engineering. Eftersom ML utgör en central del av många AI-baserade system påverkas AI 
indirekt av samma begränsningar som ML-modeller uppvisar, såsom beroende av träningsdata, modellarkitektur och kontinuerliga uppdateringar \cite{basnet2023anti_phishing, algarni2024comprehensive_taxonomy}. Om träningsdata är bristfälliga eller utsätts för en fientlig attack, kan modellen utveckla brister i 
sin klassificeringsförmåga. Därtill kvarstår det centrala problemet med AI, vilket är att hotaktörerna också använder sig av AI. Användningen av AI 
skapar då indirekt en osynlig teknologisk kapprustning där både angripare och försvarare använder sig av samma verktyg för att ligga steget före den andra \cite{zhang2024digital_deception}.

Dessutom förekommer det fler begränsningar med AI i form av dess analyseringsprocesser. AI analyserar tekniska indikatorer, medan social engineering i grunden är ett socio-tekniskt fenomen \cite{washo2021interdisciplinary}. 
Manipulationer sker i mänskliga sammanhang, organisatoriska rutiner och beslutsprocesser, där även de mest avancerade AI-systemen inte fullt kan ersätta mänskligt beslutsfattande eller  säkerhetskultur. Därpå kan nya risker 
uppstå i organisationer på grund av ett överdrivet beroende av AI. Överdrivet beroende kan leda till utvecklingen av falsk trygghet och förlitande på AI-baserat skydd, vilket därmed kan leda till minskad vaksamhet eller riskbedömning.  

\section{Autentisering och åtkomstkontroll}
Eftersom social engineering är en kombination av både psykologisk manipulation och användning av tekniska verktyg, krävs det mer avancerade lösningar för att uppnå optimalt skydd. Tekniska system, som AI-baserat skydd, 
kan identifiera och blockera många sorters angrepp, men den mänskliga faktorn kvarstår i beslutsprocesser och beteendemönster. Specifikt när det gäller organisationer spelar säkerhetskultur, utbildningsnivå och interna rutiner en avgörande 
roll i hur motståndskraftiga de är mot social engineering \cite{washo2021interdisciplinary}. Trots de olika motåtgärderna är utbildnings- och medvetandeprogram riktade till anställda bland de mest använda säkerhetsåtgärdererna. Enligt \cite{albladi2019training_awareness} är 
dessa utbildningsprogram dock bristfälliga, eftersom utbildningen ofta erbjuds endast några gånger istället för kontinuerligt. För att uppnå ett effektivt skydd bör utbildningen vara återkommande, scenariobaserad och anpassad 
efter organisationens behov. 

ENISA betonar betydelsen av att utveckla en stark cybersäkerhetskultur där säkerhetsbeteenden integreras i vardagliga rutiner \cite{enisa2018cybersecurity_culture}. Integrationen av cybersäkerhetskultur innebär att säkerheten 
inte endast ska fungera som en teknisk funktion utan även som en del av organisationen och dess ledarskap. När dessa säkerhetsprinciper tydligt är anpassade till organisationens struktur minskar sannolikheten för lyckad manipulation. 
Ett exempel på detta är användningen av simulerade nätfiskekampanjer för att mäta och förbättra arbetarnas motståndskraft \cite{albladi2018defense_mechanisms}. Genom kontrollerade tester kan organisationer identifiera sårbarheter och 
anpassa utbildningsinsatser baserat på resultaten. Samtidigt måste även dessa simuleringar genomföras med hänsyn till organisatorisk tillit och arbetsmiljö, eftersom övervakning kan påverka arbetarnas upplevelse och trygghet.

Utöver olika utbildningar är tydliga policyer och rutiner viktiga för säkerheten \cite{gururaj2024social_engineering}. Principer som lägsta behörighet, tydliga verifieringsprocesser vid ekonomiska transaktioner och 
användning av flerfaktorautentisering bidrar också till förbättrat skydd \cite{algarni2024comprehensive_taxonomy}. Genom att implementera dessa strukturella säkerhetsåtgärder kan organisationer begränsa skadan även om en 
individ blir utsatt för manipulation.